% Canonical source-aligned PageRank/RPPR reference formulation.
%
% This fragment is intentionally heading-free at the section level so that the
% active paper and every standalone note can input it. It is a manuscript
% reference convention, not an implementation-wide residual decision.
% Primary source: Fountoulakis and Martinez-Rubio, arXiv:2602.21138v2,
% PDF p. 3, Section 3 and equation (RPPR).

\subsection{Common source-aligned PageRank and RPPR model}
\label{subsec:shared-source-aligned-problem}

All documents in this manuscript workspace use the following reference
language. It follows the plain italic notation of the comparison paper:
\(x,s,A,D,Q,I\) denote column vectors or matrices as determined by their
displayed domains. This is the scoped exception to the project's usual
bold-vector and bold-matrix typography. A note may impose a stronger
assumption after this block, but it must not redefine these objects.

Let \(G=(V,E)\) be a finite, simple, undirected, unweighted graph without
isolated vertices, let \(V=[n]:=\{1,\ldots,n\}\), and let
\(A\in\{0,1\}^{n\times n}\) be its symmetric adjacency matrix. For
\(i\in V\), write
\[
    \mathcal N(i):=\{j\in V:\{i,j\}\in E\},
    \qquad
    d_i:=|\mathcal N(i)|,
    \qquad
    D:=\diag(d_1,\ldots,d_n).
\]
For \(S\subseteq V\), define
\[
    \mathcal N(S):=\bigcup_{i\in S}\mathcal N(i),
    \qquad
    \partial S:=\mathcal N(S)\setminus S,
    \qquad
    \operatorname{Ext}(S):=V\setminus(S\cup\partial S),
\]
\begin{equation}
    \vol(S):=\sum_{i\in S}d_i.
    \label{eq:shared-volume}
\end{equation}
The support of a vector is
\(\supp(x):=\{i\in[n]:x_i\neq0\}\). Matrix inequalities use the Loewner
order, and all unqualified vector inequalities are coordinatewise.

The seed is a column distribution
\begin{equation}
    s\in\R_+^n,
    \qquad
    \inner{\one}{s}=1.
    \label{eq:shared-seed}
\end{equation}
The single-seed specialization is always written \(s=e_v\), where \(v\in V\)
is the seed vertex. Thus \(s\) is never used as a vertex label. For
\(\alpha\in(0,1]\), define
\begin{equation}
\begin{aligned}
    \mathcal L&:=I-D^{-1/2}AD^{-1/2},\\
    Q&:=\alpha I+\frac{1-\alpha}{2}\mathcal L
      =\frac{1+\alpha}{2}I
       -\frac{1-\alpha}{2}D^{-1/2}AD^{-1/2},\\
    b&:=\alpha D^{-1/2}s.
\end{aligned}
    \label{eq:shared-pagerank-matrices}
\end{equation}
Since \(0\preceq\mathcal L\preceq2I\),
\(\alpha I\preceq Q\preceq I\). The shared smooth PageRank quadratic is
\begin{equation}
    f(x):=\frac12\inner{x}{Qx}-\inner{b}{x},
    \qquad
    \nabla f(x)=Qx-b.
    \label{eq:shared-pagerank-objective}
\end{equation}
Its unique unregularized minimizer and associated lazy personalized PageRank
vector are
\begin{equation}
    x^0:=Q^{-1}b,
    \qquad
    \pi:=D^{1/2}x^0.
    \label{eq:shared-ppr-solution}
\end{equation}

For a regularization parameter \(\rho>0\), reserve \(g_\rho\) for the
nonsmooth weighted \(\ell_1\) term and define
\begin{equation}
    g_\rho(x):=\alpha\rho\norm{D^{1/2}x}_1,
    \qquad
    F_\rho(x):=f(x)+g_\rho(x),
    \qquad
    x^\star(\rho):=\argmin_{x\in\R^n}F_\rho(x).
    \label{eq:shared-rppr-objective}
\end{equation}
When \(\rho\) is fixed, \(x^\star\) abbreviates \(x^\star(\rho)\), never the
unregularized point \(x^0\). Write
\[
    S^\star(\rho):=\supp(x^\star(\rho)),
    \qquad
    I^\star(\rho):=[n]\setminus S^\star(\rho).
\]
The source records \(x^\star(\rho)\geq0\),
\(\vol(S^\star(\rho))\leq1/\rho\), and the coordinatewise conditions
\begin{equation}
    \nabla_i f(x^\star(\rho))
    \in
    \begin{cases}
        \{-\alpha\rho\sqrt{d_i}\}, & x_i^\star(\rho)>0,\\
        [-\alpha\rho\sqrt{d_i},0], & x_i^\star(\rho)=0.
    \end{cases}
    \label{eq:shared-rppr-kkt}
\end{equation}

The common computational unit is one repeated adjacency-list scan: scanning
coordinate \(i\) costs \(d_i\), and scanning a set \(S\) costs \(\vol(S)\).
Iteration count and wall-clock time do not replace this degree-weighted work
measure.

Finally, accuracy symbols remain distinct. The APPR threshold is
\(\epsappr\); \(\epsobj\) denotes an objective-gap target;
\(\epspg\) denotes the source experiment's proximal fixed-point tolerance; and
\(\epsppr\) may denote a document-scoped degree-normalized PPR error target.
No equality or conversion among these quantities is assumed. In particular,
this shared model does not resolve the repository-wide residual decision.


\paragraph{Note specialization.}
This note assumes \(0<\alpha<1/2\). Define the proof-scoped degree vector and
weighted mass functional
\[
  w:=D^{1/2}\one,
  \qquad
  p(x):=w^Tx.
\]
The identities
\begin{equation}
  Qw=\alpha w,
  \qquad
  w^TQ=\alpha w^T,
  \qquad
  w^Tb=\alpha
  \label{eq:weighted-identities}
\end{equation}
are used throughout.  The matrix $Q$ is a symmetric positive-definite
$M$-matrix and satisfies $\alpha I\preceq Q\preceq I$.

\subsection{Orthant form and note-scoped KKT residual}

Write \(J(x):=\norm{D^{1/2}x}_1=\sum_i w_i|x_i|\). Since \(b\ge0\)
and \(Q\) is an \(M\)-matrix, the shared minimizer \(x^\star(\rho)\) is
nonnegative. In this note it is convenient to write it as \(x_\rho^\star\)
and use the equivalent orthant formulation
\begin{equation}
 F_\rho^+(x)
 :=\frac12x^TQx-b^Tx+\alpha\rho w^Tx+I_{\R_+^V}(x).
 \label{eq:orthant-rppr}
\end{equation}

On the orthant, introduce the note-scoped signed KKT residual
\begin{equation}
  \kappa_\rho(x):=b-Qx-\alpha\rho w.
  \label{eq:kkt-residual}
\end{equation}
The minimizer obeys
\begin{equation}
  \kappa_\rho(x_\rho^\star)\le0,
  \qquad
  x_\rho^\star\ge0,
  \qquad
  x_{\rho,i}^\star \kappa_{\rho,i}(x_\rho^\star)=0.
  \label{eq:kkt}
\end{equation}
Our sign convention makes positive residual a legal direction in which to
increase a coordinate.  This convention is local to this note.

Multiplying the active KKT equations by $w$ gives the standard support budget
\begin{equation}
  p(x_\rho^\star)+\rho\vol(\supp x_\rho^\star)\le1.
  \label{eq:rppr-support-budget}
\end{equation}
Thus $\rho^{-1}$ is simultaneously a regularization scale and a worst-case
support-volume scale. It is not identified here with an unresolved global
accuracy symbol.

\subsection{Local work}

Scanning the neighborhood of vertex $i$ costs $d_i$. By the shared work
definition, a batched active set $S$ therefore costs $\vol(S)$.
For an active-set trajectory $(S_t)_{t<T}$ we charge
\begin{equation}
  \Work_T:=\sum_{t=0}^{T-1}\vol(S_t).
  \label{eq:work-model}
\end{equation}
This is the same edge-operation unit used in the locally evolving-set
framework of \citet{zhou2024iterative}, although different algorithms may
choose different active sets.

The target~\eqref{eq:target-work} is stronger than multiplying an
$O(1/\sqrt\alpha)$ iteration count by the trivial $O(1/\rho)$ support bound
only if each accelerated subproblem can actually be solved with local work and
the momentum state itself does not create expensive transient support.
